\documentclass[10pt,a4paper]{amsart}
\usepackage[margin=19mm]{geometry}
\usepackage{amsmath,amssymb,mathtools}
\usepackage[scr=rsfs,cal=euler]{mathalfa}
\usepackage{xfrac}
\usepackage[unicode,hidelinks,bookmarks=false]{hyperref}
\newcommand{\ie}{i.e.\ }
\newcommand{\cf}{cf.\ }
\newcommand{\resp}{respectively\ }
\newcommand{\Subsection}{\subsection}
\providecommand{\nobreakdash}{\nobreak\hbox{-}}
\setlength{\emergencystretch}{2em}
\multlinegap0pt
\usepackage{lmodern}
\usepackage{mathtools}
\usepackage{bm}
\usepackage{booktabs}
\usepackage{placeins}
\hypersetup{colorlinks=true,linkcolor=blue,citecolor=blue,urlcolor=blue}
\providecommand{\Description}[1]{}
\providecommand{\textsubscript}[1]{$_{\mathrm{#1}}$}
\newtheorem{theorem}{Theorem}
\newtheorem{proposition}[theorem]{Proposition}
\newtheorem{corollary}[theorem]{Corollary}
\newtheorem{definition}[theorem]{Definition}
\newtheorem{remark}[theorem]{Remark}
\newcommand{\SPINE}{\textsc{SPINE}}
\newcommand{\Neff}{N_{\mathrm{eff}}}
\newcommand{\aeff}{a_{\mathrm{eff}}}
\newcommand{\rhostar}{\rho^{\star}}
\newcommand{\cstar}{c^{\star}}
\newcommand{\cO}{\mathcal O}
\newcommand{\cC}{\mathcal C}
\newcommand{\cB}{\mathcal B}
\newcommand{\Rplus}{\mathbb R_{\ge 0}}
\begin{document}
\title[Scale Partitioning by Incremental Nested Entropy: A Measure-]{Scale Partitioning by Incremental Nested Entropy: A Measure-Oriented Theory of Multiscale Structure}
\author{Abd AlRahman R. AlMomani}
\date{}
\maketitle
\noindent\emph{Source and licence.} Scale Partitioning by Incremental Nested Entropy: A Measure-Oriented Theory of Multiscale Structure. Version arXiv:2608.17391v1 (CC BY 4.0). This material is distributed under the Creative Commons Attribution 4.0 International licence, http://creativecommons.org/licenses/by/4.0/, as recorded in the arXiv OAI licence metadata for this exact version and shown on the abstract page https://arxiv.org/abs/2608.17391v1. Full rights notice: licenses/arxiv-2608.17391-CC-BY-4.0.txt.\par\smallskip
\noindent\textit{Editorial scope.} Counted unit: the original \texttt{theorem} environment \texttt{thm:recurrence} at source lines 118--130 together with its complete proof at lines 131--133; the delivered document keeps source lines 87--134 so that every referenced label and every citation resolves inside it. Native editable source is the paper's own concatenated TeX; the AMS wrapper is editorial formatting only. No publisher class, upstream script or unrelated artwork is redistributed.
\par\smallskip\noindent\textit{Reference note.} This article ships no compiled bibliography (biblatex); the entries delivered here are composed from the author's own BibTeX (.bib) fields, with no field invented and no entry dropped.
\section{Nested entropy of an ordered measure and the entropy-effective state}
\label{sec:nested}

Let \(\cO_N=\{o_1,\ldots,o_N\}\) be a finite set of objects, and let \(\mu:\cO_N\to\Rplus\) be a declared scientific measure, with \(m_i=\mu(o_i)\). The map \(\mu\) assigns a nonnegative weight to each object; it is not a probability distribution. The measure may carry physical units and may arise from geometry, dynamics, counts, spectral decomposition, network structure, or another application-specific construction. Probability enters only after normalization of a prefix. Let \(m_{(1)}\le\cdots\le m_{(N)}\) denote a stable nondecreasing ordering of the measure values, with object identities carried through the ordering, and let \(S_n=\sum_{j=1}^n m_{(j)}\) be the cumulative measure of the first \(n\) ordered objects. % Orientation does not alter this ordering; it is applied after partitioning to interpret the discovered strata.

Whenever \(S_n>0\), the first \(n\) ordered values induce normalized weights \(p_j^{(n)}=m_{(j)}/S_n\). These weights sum to one and therefore define a finite probability vector on the prefix, but their meaning is purely measure-normalized. Unless \(\mu\) is itself an occurrence frequency, \(p_j^{(n)}\) should not be read as the probability that object \(o_{(j)}\) occurs. Rather, it is the fraction of accumulated prefix measure carried by that object. This distinction is useful in applications because the same entropy construction can act on measures with very different scientific meanings.

\begin{definition}[Incremental nested entropy]
\label{def:nested}
The nested entropy profile is the sequence \(E_1,\ldots,E_N\) obtained by evaluating Shannon entropy on each normalized ordered prefix,
\begin{equation}
E_n=-\sum_{j=1}^{n}p_j^{(n)}\log_2 p_j^{(n)}.
\label{eq:nested-entropy}
\end{equation}
If \(S_n=0\), we set \(E_n=0\). Zero-mass terms contribute zero by continuity. For \(n\ge2\), the local entropy increment is \(\Delta_n=E_n-E_{n-1}\).
\end{definition}

The profile in Eq.~\eqref{eq:nested-entropy} is a sequence of entropies of different normalized allocations, not an entropy of rank and not a histogram entropy of the original sample. Ordering determines which object enters next, while renormalization determines how the accumulated measure is redistributed after that entry. If all \(n\) positive values in a prefix are equal, then the normalized allocation is uniform and \(E_n=\log_2 n\). An exactly homogeneous measure scale therefore produces a strictly increasing profile. At the opposite extreme, an incoming value that captures nearly all of the enlarged prefix measure drives the normalized allocation toward a single dominant component and can reduce entropy sharply. The sign of each local entropy increment is determined by the competition between these two tendencies.

This competition is explicit in the one-step recurrence, which is the Shannon grouping identity specialized to the ascending nested-prefix construction \cite{shannon1948,coverthomas2006}. Let \(q_n=m_{(n)}/S_n\) be the normalized mass of the incoming value after it has entered the prefix, and let \(h_2(q)=-q\log_2q-(1-q)\log_2(1-q)\) denote binary entropy.

% \begin{theorem}[One-step entropy recurrence]
% \label{thm:recurrence}
% For \(n\ge2\) with \(S_n>0\),
% \begin{equation}
% E_n=(1-q_n)E_{n-1}+h_2(q_n),
% \qquad
% \Delta_n=h_2(q_n)-q_nE_{n-1}.
% \label{eq:recurrence}
% \end{equation}
% \end{theorem}

\begin{theorem}[One-step entropy recurrence]
\label{thm:recurrence}
For \(n\ge2\) with \(S_n>0\),
\begin{equation}
\begin{aligned}
E_n
&=(1-q_n)E_{n-1}+h_2(q_n),\\
\Delta_n
&=h_2(q_n)-q_nE_{n-1}.
\end{aligned}
\label{eq:recurrence}
\end{equation}
\end{theorem}

\begin{proof}
If \(S_{n-1}=0\), then \(q_n=1\) and both identities hold with \(E_{n-1}=E_n=h_2(1)=0\). Otherwise, after the new value enters, each of the first \(n-1\) normalized masses is multiplied by the common factor \(1-q_n\), while the new component has normalized mass \(q_n\). Expanding the Shannon entropy of the vector \(((1-q_n)\bm p^{(n-1)},q_n)\) gives the first identity. Subtracting \(E_{n-1}\) gives the second. A closed-form identity equivalent to the recurrence is given in Supplementary Note~S1.1.
\end{proof}


\begin{thebibliography}{99}
\bibitem[]{coverthomas2006} Cover, Thomas M.; Thomas, Joy A.. Elements of Information Theory. (2006)
\bibitem[]{shannon1948} Shannon, Claude E.. A Mathematical Theory of Communication. Bell System Technical Journal 27 (1948) 379-423
\end{thebibliography}
\end{document}
